\documentclass[letterpaper, twoside]{article}

\usepackage{graphicx}
\usepackage[margin=2.54cm]{geometry}
\usepackage{mathtools}
\usepackage{amssymb}
\usepackage{amsthm}
\usepackage{fancyhdr}
\usepackage{setspace}
\usepackage{parskip}
\usepackage{caption}


%Encabezado
\pagestyle{fancy}
\fancyhf{}
\lhead{Actividad 1}
\rhead{Material 1}
\cfoot{\thepage}

%Margenes de las tablas
\setlength{\tabcolsep}{10pt} % Aumenta el margen horizontal
\renewcommand{\arraystretch}{1.7} % Aumenta la altura de las filas


\begin{document}

%.................Material 1.................................

\begin{center}
  \textbf{\large Material 1}
\end{center}

\vspace{.5cm}

Objetivo: que el alumno entienda, a grandes rasgos, cómo funciona la tabla de verdad de la negación.

\vspace{.5cm}

\textbf{Proposición}

Una proposición es un enunciado del cual se puede decir si es verdadera o falsa. Una forma muy simple para mencionar a las proposiciones sin tener que estarlas escribiendo todo el tiempo es usando las letras mayúsculas $P,Q,R,S$.\par

\vspace{.5cm}

Como ya se menciono, una proposición $P$ puede ser verdadera ($V$) o falsa ($F$), es decir, tiene dos valores de verdad. Se puede representar con la siguiente tabla:\par

\begin{table}[h]
  \centering
    \begin{tabular}{|c|}
      \hline
      $P$ \\ \hline
      $V$ \\ \hline
      $F$ \\
      \hline
    \end{tabular}
  \caption*{Tabla de verdad de $P$}
\end{table}

\vspace{1cm}

\textit{Ejercicio} 1: Completa las tablas de verdad de las proposiciones $Q$, $R$ y $S$ colocando el valor de verdad faltante.

\begin{table}[htbp]

  \centering
    \begin{minipage}[t]{0.31\textwidth}
      \centering
      \begin{tabular}{|c|}
        \hline
        $Q$ \\ \hline
        $V$ \\ \hline
            \\ 
        \hline
      \end{tabular}
    \end{minipage}%
    \hfill
    \begin{minipage}[t]{0.31\textwidth}
      \centering
      \begin{tabular}{|c|}
        \hline
        $R$ \\ \hline
        $F$ \\ \hline
            \\ 
        \hline
      \end{tabular}
    \end{minipage}%
    \hfill
    \begin{minipage}[t]{0.31\textwidth}
      \centering
      \begin{tabular}{|c|}
        \hline
        $S$ \\ \hline
            \\ \hline
            \\ 
        \hline
      \end{tabular}
    \end{minipage}

\end{table}

\vspace{1cm}

\textbf{Tabla de verdad de la Negación}

La negación de una proposición es una nueva proposición.\par

\vspace{.5cm}

El símbolo $\boldsymbol{\neg}$ representa la negación y es llamado "conectivo de la negación". De está manera, para una proposición $P$, la forma de representar su negación es de la siguiente manera:\par
\begin{center}
  $\neg P$
\end{center}

Se lee: no $P$

\vspace{.5cm}

Los valores de verdad de la proposición $\neg P$ se resumen en los siguientes dos casos:
\begin{itemize}
  \item Caso 1: Si la proposición $P$ es verdadera entonces su negación, $\neg P$, resulta ser falsa.\par
  \begin{center}
    \begin{tabular}{|c|c|}
    \hline
    $P$ & $\neg P$ \\ \hline
    $V$ & $F$ \\
    \hline
  \end{tabular}
  \end{center}
  \item Caso 2: Si la proposición $P$ es falsa entonces su negación, $\neg P$, resulta ser verdadera.\par
  \begin{center}
    \begin{tabular}{|c|c|}
    \hline
    $P$ & $\neg P$ \\ \hline
    $F$ & $V$ \\
    \hline
  \end{tabular}
  \end{center}
\end{itemize}

\vspace{.5cm}

Esto se puede visualizar de mejor manera utilizando la tabla de verdad de la negación.

\begin{table}[h]
  \centering
    \begin{tabular}{|c|c|c|}
      \hline
      Casos  & $P$ & $\neg P$ \\ \hline
      Caso 1 & $V$ & $F$      \\ \hline
      Caso 2 & $F$ & $V$       \\
      \hline
    \end{tabular}
  \caption*{Tabla de verdad de la negación $\neg P$}
\end{table}

\vspace{.5cm}

Observa que la tabla de arriba esta compuesta por columnas; una para la proposición $P$ y otra para la proposición $\neg P$; y filas con los valores de verdad posibles de $P$ y $\neg P$. La columna de "Casos" solo es de apoyo para poder relacionar de mejor manera el texto con la fila correspondiente.

\vspace{.5cm}

\textit{Ejercicio} 2: Apóyate en la tabla de verdad de la negación y responde:\par

\begin{enumerate}
  \item Si la proposición $\neg P$ es verdadera entonces la proposición $P$ es \hrulefill
  \item Si la proposición $\neg P$ es falsa entonces la proposición $P$ es \hrulefill
\end{enumerate}

\vspace{.5cm}

\textit{Ejercicio} 3: Completa las tablas de verdad con los valores de verdad faltantes.

\begin{table}[h]
  
  \centering
  \begin{minipage}[t]{0.31\textwidth}
    \centering
    \begin{tabular}{|c|c|}
      \hline
      $Q$ & $\neg Q$ \\ \hline
      $V$ & \\ \hline
      $F$ & \\
      \hline
    \end{tabular}
  \end{minipage}
  \hfil
  \begin{minipage}[t]{0.31\textwidth}
    \centering
    \begin{tabular}{|c|c|}
      \hline
      $R$ & $\neg R$ \\ \hline
          & $F$ \\ \hline
      $F$ & \\
      \hline
    \end{tabular}
  \end{minipage}
  \begin{minipage}[t]{0.31\textwidth}
    \centering
    \begin{tabular}{|c|c|}
      \hline
      $\neg S$ & $S$ \\ \hline
        & $V$\\ \hline
        & $F$ \\
      \hline
    \end{tabular}
  \end{minipage}

\end{table}

%.................................Material 2................................

\newpage

\clearpage
\fancyhf{}
\lhead{Actividad 1}
\rhead{Material 2}
\setcounter{page}{1} %Reinicio del contador de página
\cfoot{\thepage}

\begin{center}
  \textbf{\large Material 2}
\end{center}

\vspace{.5cm}

Objetivo: que el alumno entienda, a grandes rasgos, cómo funciona la tabla de verdad de la disyunción.

\vspace{.5cm}

\textbf{Proposición}

Una proposición es un enunciado del cual se puede decir si es verdadera o falsa. Una forma muy simple para mencionar a las proposiciones sin tener que estarlas escribiendo todo el tiempo es usando las letras mayúsculas $P,Q,R,S$.\par

\vspace{.5cm}

Como ya se menciono, una proposición $P$ puede ser verdadera ($V$) o falsa ($F$), es decir, tiene dos valores de verdad. Se puede representar en la siguiente tabla:\par

\begin{table}[h]
  \centering
    \begin{tabular}{|c|}
      \hline
      $P$ \\ \hline
      $V$ \\ \hline
      $F$ \\
      \hline
    \end{tabular}
  \caption*{Tabla de verdad de $P$}
\end{table}

\vspace{1cm}

\textit{Ejercicio} 1: Completa las tablas de verdad de las proposiciones $Q$, $R$ y $S$ colocando el valor de verdad faltante.

\begin{table}[htbp]

  \centering
    \begin{minipage}[t]{0.31\textwidth}
      \centering
      \begin{tabular}{|c|}
        \hline
        $Q$ \\ \hline
        $V$ \\ \hline
            \\ 
        \hline
      \end{tabular}
    \end{minipage}%
    \hfill
    \begin{minipage}[t]{0.31\textwidth}
      \centering
      \begin{tabular}{|c|}
        \hline
        $R$ \\ \hline
        $F$ \\ \hline
            \\ 
        \hline
      \end{tabular}
    \end{minipage}%
    \hfill
    \begin{minipage}[t]{0.31\textwidth}
      \centering
      \begin{tabular}{|c|}
        \hline
        $S$ \\ \hline
            \\ \hline
            \\ 
        \hline
      \end{tabular}
    \end{minipage}

\end{table}

\vspace{1cm}

\textbf{ Conectivo de la Disyunción}\par

\vspace{.5cm}

Al considerar a la proposición $P$ y a la proposición $Q$ se puede formar una nueva proposición utilizando el conectivo de la disyunción.\par

\vspace{.5cm}

El símbolo $\boldsymbol{\lor}$ representa al conectivo de la disyunción. De está manera, para una proposición $P$ y otra proposición $Q$, la forma de representar la disyunción entre estas dos proposiciones es la siguiente:\par

\begin{center}
  $P\lor Q$
\end{center}

Se lee: $P$ ó $Q$

\vspace{.5cm}

\newpage

Los valores de verdad de la proposición $P\lor Q$ depende de los valores de verdad de cada una de las proposiciones que la componen y de las combinaciones de estos valores de verdad. Así, se tienen los siguientes 4 casos:

\begin{itemize}
  \item Caso 1: Si las proposiciones $P$ y $Q$ son verdaderas entonces resulta que la proposición $P\lor Q$ es verdadera.
  \begin{center}
    \begin{tabular}{|c|c|c|}
    \hline
    $P$ & $Q$ & $P\lor Q$ \\ \hline
    $V$ & $V$ & $V$\\
    \hline
  \end{tabular}
  \end{center}
\vspace{.7cm}
  \item Caso 2: Si la proposición $P$ es verdadera y la proposición $Q$ es falsa entonces la proposición $P\lor Q$ es verdadera.
  \begin{center}
    \begin{tabular}{|c|c|c|}
    \hline
    $P$ & $Q$ & $P\lor Q$ \\ \hline
    $V$ & $F$ & $V$\\
    \hline
  \end{tabular}
  \end{center}
\vspace{.7cm}
  \item Caso 3: Si la proposición $P$ es falsa y la proposición $Q$ es verdadera entonces la proposición $P\lor Q$ es verdadera.
  \begin{center}
    \begin{tabular}{|c|c|c|}
    \hline
    $P$ & $Q$ & $P\lor Q$ \\ \hline
    $F$ & $V$ & $V$\\
    \hline
  \end{tabular}
  \end{center}
\vspace{.7cm}
  \item Caso 4: Si las proposiciones $P$ y $Q$ son falsas entonces resulta que la proposición $P\lor Q$ es falsa.
  \begin{center}
    \begin{tabular}{|c|c|c|}
    \hline
    $P$ & $Q$ & $P\lor Q$ \\ \hline
    $F$ & $F$ & $F$\\
    \hline
  \end{tabular}
  \end{center}

\end{itemize}

\vspace{.5cm}

Esto se puede visualizar de mejor manera utilizando la tabla de verdad de la disyunción.

\begin{table}[h]
  \centering
  \begin{tabular}{|c|c|c|c|}
    \hline
    Casos  & $P$ & $Q$ & $P\lor Q$ \\ \hline
    Caso 1 & $V$ & $V$ & $V$ \\ \hline
    Caso 2 & $V$ & $F$ & $V$ \\ \hline
    Caso 3 & $F$ & $V$ & $V$ \\ \hline
    Caso 4 & $F$ & $F$ & $F$ \\ \hline 
  \end{tabular}
  \caption*{Tabla de verdad de la disyunción $P\lor Q$}
\end{table}


Observa que la tabla de arriba esta compuesta por columnas; una columna para $P$, otra para $Q$ y una más para $P\lor Q$; y filas para los valores de verdad posibles de $P$, $Q$ y $P\lor Q$. La columna de "Casos" solo es de apoyo para poder relacionar de mejor manera el texto con la fila correspondiente.

\newpage

\textit{Ejercicio} 2: Apóyate en la tabla de verdad de la disyunción y responde:

\begin{enumerate}
  \item ¿En qué casos la proposición $P\lor Q$ es verdadera?
\vspace{5cm}
  \item ¿En qué casos la proposición $P\lor Q$ es falsa?
\end{enumerate}

\vspace{5cm}

\textit{Ejercicio} 3: Completa las tablas de verdad con los valores de verdad faltantes.

\begin{table}[h]
  
  \centering
  \begin{minipage}[t]{0.31\textwidth}
    \centering
    \begin{tabular}{|c|c|c|}
      \hline
      $P$ & $Q$ & $P\lor Q$ \\ \hline
      $V$ & $V$ &  \\ \hline
      $V$ & $F$ &  \\ \hline
      $F$ & $V$ &  \\ \hline
      $F$ & $F$ &  \\
      \hline
    \end{tabular}
  \end{minipage}
  \hfil
  \begin{minipage}[t]{0.31\textwidth}
    \centering
    \begin{tabular}{|c|c|c|}
      \hline
      $Q$ & $S$ & $Q\lor S$ \\ \hline
      $F$ & $F$ &  \\ \hline
      $F$ & $V$ &  \\ \hline
      $V$ & $F$ &  \\ \hline
      $V$ & $V$ &  \\
      \hline
    \end{tabular}
  \end{minipage}
  \begin{minipage}[t]{0.31\textwidth}
    \centering
    \begin{tabular}{|c|c|c|}
      \hline
      $Q$ & $S$ & $Q\lor S$ \\ \hline
      $F$ &     &  $V$ \\ \hline
          & $F$ &  $F$ \\ \hline
          & $F$ &  $V$ \\ \hline
      $V$ & $V$ &  \\
      \hline
    \end{tabular}
  \end{minipage}

\end{table}
%..................................Material 3........................................

\newpage

\clearpage
\fancyhf{}
\lhead{Actividad 1}
\rhead{Material 3}
\setcounter{page}{1} %Reinicio del contador de página
\cfoot{\thepage}

\begin{center}
  \textbf{\large Material 3}
\end{center}

\vspace{.5cm}

Objetivo: que el alumno entienda, a grandes rasgos, cómo funciona la tabla de verdad de la disyunción.

\vspace{.5cm}

\textbf{Proposición}

Una proposición es un enunciado del cual se puede decir si es verdadera o falsa. Una forma muy simple para mencionar a las proposiciones sin tener que estarlas escribiendo todo el tiempo es usando las letras mayúsculas $P,Q,R,S$.\par

\vspace{.5cm}

Como ya se menciono, una proposición $P$ puede ser verdadera ($V$) o falsa ($F$), es decir, tiene dos valores de verdad. Se puede representar en la siguiente tabla:\par

\begin{table}[h]
  \centering
    \begin{tabular}{|c|}
      \hline
      $P$ \\ \hline
      $V$ \\ \hline
      $F$ \\
      \hline
    \end{tabular}
  \caption*{Tabla de verdad de $P$}
\end{table}

\vspace{1cm}

\textit{Ejercicio} 1: Completa las tablas de verdad de las proposiciones $Q$, $R$ y $S$ colocando el valor de verdad faltante.

\begin{table}[htbp]

  \centering
    \begin{minipage}[t]{0.31\textwidth}
      \centering
      \begin{tabular}{|c|}
        \hline
        $Q$ \\ \hline
        $V$ \\ \hline
            \\ 
        \hline
      \end{tabular}
    \end{minipage}%
    \hfill
    \begin{minipage}[t]{0.31\textwidth}
      \centering
      \begin{tabular}{|c|}
        \hline
        $R$ \\ \hline
        $F$ \\ \hline
            \\ 
        \hline
      \end{tabular}
    \end{minipage}%
    \hfill
    \begin{minipage}[t]{0.31\textwidth}
      \centering
      \begin{tabular}{|c|}
        \hline
        $S$ \\ \hline
            \\ \hline
            \\ 
        \hline
      \end{tabular}
    \end{minipage}

\end{table}

\vspace{1cm}

\textbf{ Conectivo de la Conjunción}\par

\vspace{.5cm}

Al considerar a la proposición $P$ y a la proposición $Q$ se puede formar una nueva proposición utilizando el conectivo de la conjunción.\par

\vspace{.5cm}

El símbolo $\boldsymbol{\land}$ representa al conectivo de la conjunción. De está manera, para una proposición $P$ y otra proposición $Q$, la forma de representar la conjunción entre estas dos proposiciones es la siguiente:\par

\begin{center}
  $P\land Q$
\end{center}

Se lee: $P$ y $Q$

\vspace{.5cm}

\newpage

Los valores de verdad de la proposición $P\land Q$ depende de los valores de verdad de cada una de las proposiciones que la componen y de las combinaciones de estos valores de verdad. Así, se tienen los siguientes 4 casos:

\begin{itemize}
  \item Caso 1: Si las proposiciones $P$ y $Q$ son verdaderas entonces resulta que la proposición $P\land Q$ es verdadera.
  \begin{center}
    \begin{tabular}{|c|c|c|}
    \hline
    $P$ & $Q$ & $P\land Q$ \\ \hline
    $V$ & $V$ & $V$\\
    \hline
  \end{tabular}
  \end{center}
\vspace{.7cm}
  \item Caso 2: Si la proposición $P$ es verdadera y la proposición $Q$ es falsa entonces la proposición $P\land Q$ es falsa.
  \begin{center}
    \begin{tabular}{|c|c|c|}
    \hline
    $P$ & $Q$ & $P\land Q$ \\ \hline
    $V$ & $F$ & $F$\\
    \hline
  \end{tabular}
  \end{center}
\vspace{.7cm}
  \item Caso 3: Si la proposición $P$ es falsa y la proposición $Q$ es verdadera entonces la proposición $P\lor Q$ es falsa.
  \begin{center}
    \begin{tabular}{|c|c|c|}
    \hline
    $P$ & $Q$ & $P\land Q$ \\ \hline
    $F$ & $V$ & $F$\\
    \hline
  \end{tabular}
  \end{center}
\vspace{.7cm}
  \item Caso 4: Si las proposiciones $P$ y $Q$ son falsas entonces resulta que la proposición $P\lor Q$ es falsa.
  \begin{center}
    \begin{tabular}{|c|c|c|}
    \hline
    $P$ & $Q$ & $P\land Q$ \\ \hline
    $F$ & $F$ & $F$\\
    \hline
  \end{tabular}
  \end{center}

\end{itemize}

\vspace{.5cm}

Esto se puede visualizar de mejor manera utilizando la tabla de verdad de la conjunción.

\begin{table}[h]
  \centering
  \begin{tabular}{|c|c|c|c|}
    \hline
    Casos  & $P$ & $Q$ & $P\land Q$ \\ \hline
    Caso 1 & $V$ & $V$ & $V$ \\ \hline
    Caso 2 & $V$ & $F$ & $F$ \\ \hline
    Caso 3 & $F$ & $V$ & $F$ \\ \hline
    Caso 4 & $F$ & $F$ & $F$ \\ \hline 
  \end{tabular}
  \caption*{Tabla de verdad de la disyunción $P\land Q$}
\end{table}


Observa que la tabla de arriba esta compuesta por columnas; una columna para $P$, otra para $Q$ y una más para $P\land Q$; y filas para los valores de verdad posibles de $P$, $Q$ y $P\land Q$. La columna de "Casos" solo es de apoyo para poder relacionar de mejor manera el texto con la fila correspondiente.

\newpage

\textit{Ejercicio} 2: Apóyate en la tabla de verdad de la conjunción y responde:

\begin{enumerate}
  \item ¿En qué casos la proposición $P\land Q$ es verdadera?
\vspace{5cm}
  \item ¿En qué casos la proposición $P\land Q$ es falsa?
\end{enumerate}

\vspace{5cm}

\textit{Ejercicio} 3: Completa las tablas de verdad con los valores de verdad faltantes.

\begin{table}[h]
  
  \centering
  \begin{minipage}[t]{0.31\textwidth}
    \centering
    \begin{tabular}{|c|c|c|}
      \hline
      $P$ & $Q$ & $P\land Q$ \\ \hline
      $V$ & $V$ &  \\ \hline
      $V$ & $F$ &  \\ \hline
      $F$ & $V$ &  \\ \hline
      $F$ & $F$ &  \\
      \hline
    \end{tabular}
  \end{minipage}
  \hfil
  \begin{minipage}[t]{0.31\textwidth}
    \centering
    \begin{tabular}{|c|c|c|}
      \hline
      $Q$ & $S$ & $Q\land S$ \\ \hline
      $F$ & $F$ &  \\ \hline
      $F$ & $V$ &  \\ \hline
      $V$ & $F$ &  \\ \hline
      $V$ & $V$ &  \\
      \hline
    \end{tabular}
  \end{minipage}
  \begin{minipage}[t]{0.31\textwidth}
    \centering
    \begin{tabular}{|c|c|c|}
      \hline
      $Q$ & $S$ & $Q\land S$ \\ \hline
      $V$ &     &  $F$ \\ \hline
          & $F$ &  $F$ \\ \hline
          & $V$ &  $V$ \\ \hline
      $F$ & $V$ &  \\
      \hline
    \end{tabular}
  \end{minipage}

\end{table}


%...............................Actividad 2.......................................

\newpage

\clearpage
\fancyhf{}
\lhead{Actividad 2}
%\rhead{Material 2}
\setcounter{page}{1} %Reinicio del contador de página
\cfoot{\thepage}


\textit{Ejercicio} 1: Explica a tus compañeros de equipo la tabla de verdad 






\end{document}